\chapter{INTRODUCTION}

\section{Background and Literature Review}
The predictive modeling of financial time-series has evolved from classical statistical models toward complex machine learning ensembles and natural language processing (NLP). Central to this field is the \textbf{Efficient Market Hypothesis (EMH)} \citep{fama1970}, which posits that if markets are 100\% efficient, future prices follow a random walk and are mathematically unpredictable. However, contemporary research in \textbf{Behavioral Finance} \citep{shiller2003} suggests that information asymmetry and irrational investor sentiment create exploitable "noise" in market data.

Past research has explored deep learning architectures, such as **Long Short-Term Memory (LSTM)** networks \citep{hochreiter1997}, to capture temporal dependencies. While theoretically robust for long-term memory, LSTMs often suffer from vanishing gradients and require massive datasets to achieve generalization, which is a significant **con** in the high-noise environment of daily stock prices. In contrast, **Gradient Boosting Decision Trees (GBDT)**, specifically **XGBoost** \citep{chen2016}, offer superior **pros** in handling tabular data and multi-collinearity without the "black-box" risk of over-parameterized neural networks. 

Furthermore, the integration of **Transformers** \citep{vaswani2017} has revolutionized financial NLP. While models like BERT provide state-of-the-art accuracy in sentiment extraction, their computational complexity and the need for domain-specific fine-tuning (e.g., FinBERT) remain key challenges for real-time systems. This project explores the synergy between domain-specific sentiment (FinBERT) and robust ensemble methods (XGBoost) to address the limitations of technical-only baseline models.

\section{Problem Statement and Motivation}
Traditional financial models often rely on internal market variables (Open, High, Low, Close) but fail to anticipate exogenous shocks driven by news headlines or macroeconomic shifts. This study aims to build a scientific data fusion pipeline that bridges the gap between quantitative technical indicators and qualitative textual sentiment. 

\section{Research Objectives}
To address the research problem, the following primary objectives were established:
\begin{itemize}
    \item \textbf{Data Integration}: Harmonize heterogeneous market, macro, and sentiment data streams for a multi-asset intelligence system.
    \item \textbf{Sentiment Quantification}: Utilize FinBERT to transform news headlines into numeric polarity scores to mitigate behavioral information asymmetry.
    \item \textbf{Predictive Excellence}: Scientifically validate an XGBoost ensemble classifier across five asset classes (Stocks, Gold, Crypto) using rigorous time-series cross-validation.
\end{itemize}